% GENERATED FILE. Edit canonical lessons, then run npm run generate:books.
% canonical-source-hash: fnv1a64:83287ecbdb582444
% canonical-lessons: LA-C180-commodus, LA-C180-urbanus, LA-C180-rusticus, LA-C180-hilaris, LA-C180-verecundus

\chapter{Convenient, Polite, Of the country, Cheerful, Shy}
\label{ch:la-convenient-polite-of-the-country-cheerful-shy}
\hlchaptermodality{180}

\begin{chapteropening}
\textbf{By the end of this chapter:} \emph{I can say convenient, polite, of the country, cheerful and shy.}

Complete the last lesson of chapter 180: I can say convenient, polite, of the country, cheerful and shy.
\end{chapteropening}

\section[commodus, commoda, commodum]{commodus, commoda, commodum — convenient, comfortable}

\label{lesson:LA-C180-commodus}



\begin{quote}
Before the new one: say the Latin for very, then the Latin for quite.
\end{quote}

\subsection*{You'll want to know: commodus, commoda, commodum}
\textbf{commodus, commoda, commodum} — \textquotedblleft{}convenient, comfortable\textquotedblright{}.

\begin{cousinweb}
\textbf{commodus} is \emph{com-}, \textquotedblleft{}with\textquotedblright{}, and \emph{modus}, \textquotedblleft{}measure\textquotedblright{} (Lewis \& Short): having its due measure, so fitting.
\end{cousinweb}

\begin{grammarlens}[title={What you've built}]
One more word for this chapter.
\end{grammarlens}

\subsection*{Your turn}
\begin{itemize}
\raggedright
  \item \emph{Say it:} \emph{commodus, commoda, commodum}
  \item \emph{Say it:} \emph{commodus, commoda, commodum}, once more
  \item \emph{Say it:} say \emph{commodus, commoda, commodum}
  \item \emph{Recall:} say the Latin for very, then the Latin for quite, then say \emph{commodus, commoda, commodum} again
\end{itemize}

\begin{tcolorbox}[breakable,colback=teal!4,colframe=teal!35!black,title={Before you move on}]
What does commodus, commoda, commodum mean? (\textquotedblleft{}Convenient, comfortable\textquotedblright{}.) Say it once more.
\end{tcolorbox}
\section[urbānus, urbāna, urbānum]{urbānus, urbāna, urbānum — polite, of the city}

\label{lesson:LA-C180-urbanus}



\begin{quote}
Before the new one: say the Latin for quite, then the Latin for convenient.
\end{quote}

\subsection*{You'll want to know: urbānus, urbāna, urbānum}
\textbf{urbānus, urbāna, urbānum} — \textquotedblleft{}polite, of the city\textquotedblright{}.

\begin{cousinweb}
\textbf{urbānus} is built on \emph{urbs}, \textquotedblleft{}a city\textquotedblright{} (Lewis \& Short), with \emph{-ānus}, \textquotedblleft{}belonging to\textquotedblright{}: of the city, then polished, witty. English \emph{urbane} and \emph{urban} both come from it.
\end{cousinweb}

\begin{grammarlens}[title={What you've built}]
One more word for this chapter.
\end{grammarlens}

\subsection*{Your turn}
\begin{itemize}
\raggedright
  \item \emph{Say it:} \emph{urbānus, urbāna, urbānum}
  \item \emph{Say it:} \emph{urbānus, urbāna, urbānum}, once more
  \item \emph{Say it:} say \emph{urbānus, urbāna, urbānum}
  \item \emph{Recall:} say the Latin for quite, then the Latin for convenient, then say \emph{urbānus, urbāna, urbānum} again
\end{itemize}

\begin{tcolorbox}[breakable,colback=teal!4,colframe=teal!35!black,title={Before you move on}]
What does urbānus, urbāna, urbānum mean? (\textquotedblleft{}Polite, of the city\textquotedblright{}.) Say it once more.
\end{tcolorbox}
\section[rūsticus, rūstica, rūsticum]{rūsticus, rūstica, rūsticum — of the country, boorish}

\label{lesson:LA-C180-rusticus}



\begin{quote}
Before the new one: say the Latin for convenient, then the Latin for polite.
\end{quote}

\subsection*{You'll want to know: rūsticus, rūstica, rūsticum}
\textbf{rūsticus, rūstica, rūsticum} — \textquotedblleft{}of the country, boorish\textquotedblright{}.

\begin{cousinweb}
\textbf{rūsticus} is built on \emph{rūs}, \textquotedblleft{}the countryside\textquotedblright{} (Lewis \& Short), which is usually traced to Proto-Indo-European \emph{*Hrewos}, \textquotedblleft{}open space, field\textquotedblright{} (EtymDB 2.1). English \emph{room} is a cousin.
\end{cousinweb}

\begin{grammarlens}[title={What you've built}]
One more word for this chapter.
\end{grammarlens}

\subsection*{Your turn}
\begin{itemize}
\raggedright
  \item \emph{Say it:} \emph{rūsticus, rūstica, rūsticum}
  \item \emph{Say it:} \emph{rūsticus, rūstica, rūsticum}, once more
  \item \emph{Say it:} say \emph{rūsticus, rūstica, rūsticum}
  \item \emph{Recall:} say the Latin for convenient, then the Latin for polite, then say \emph{rūsticus, rūstica, rūsticum} again
\end{itemize}

\begin{tcolorbox}[breakable,colback=teal!4,colframe=teal!35!black,title={Before you move on}]
What does rūsticus, rūstica, rūsticum mean? (\textquotedblleft{}Of the country, boorish\textquotedblright{}.) Say it once more.
\end{tcolorbox}
\section[hilaris, hilare]{hilaris, hilare — cheerful}

\label{lesson:LA-C180-hilaris}



\begin{quote}
Before the new one: say the Latin for polite, then the Latin for of the country.
\end{quote}

\subsection*{You'll want to know: hilaris, hilare}
\textbf{hilaris, hilare} — \textquotedblleft{}cheerful\textquotedblright{}.

\begin{cousinweb}
\textbf{hilaris} was borrowed from Greek \emph{hilaros}, \textquotedblleft{}cheerful\textquotedblright{} (Lewis \& Short), from \emph{hilaos}, \textquotedblleft{}gracious, kindly\textquotedblright{} (EtymDB 2.1). English \emph{hilarious} comes from it.
\end{cousinweb}

\begin{grammarlens}[title={What you've built}]
One more word for this chapter.
\end{grammarlens}

\subsection*{Your turn}
\begin{itemize}
\raggedright
  \item \emph{Say it:} \emph{hilaris, hilare}
  \item \emph{Say it:} \emph{hilaris, hilare}, once more
  \item \emph{Say it:} say \emph{hilaris, hilare}
  \item \emph{Recall:} say the Latin for polite, then the Latin for of the country, then say \emph{hilaris, hilare} again
\end{itemize}

\begin{tcolorbox}[breakable,colback=teal!4,colframe=teal!35!black,title={Before you move on}]
What does hilaris, hilare mean? (\textquotedblleft{}Cheerful\textquotedblright{}.) Say it once more.
\end{tcolorbox}
\section[verēcundus, verēcunda, verēcundum]{verēcundus, verēcunda, verēcundum — shy, modest}

\label{lesson:LA-C180-verecundus}



\begin{quote}
Before the new one: say the Latin for of the country, then the Latin for cheerful.
\end{quote}

\subsection*{You'll want to know: verēcundus, verēcunda, verēcundum}
\textbf{verēcundus, verēcunda, verēcundum} — \textquotedblleft{}shy, modest\textquotedblright{}.

\begin{cousinweb}
\textbf{verēcundus} is built on \emph{vereor}, \textquotedblleft{}to stand in awe of, to fear\textquotedblright{} (Lewis \& Short), with \emph{-cundus}, an ending for a lasting disposition. \emph{vereor} is usually traced to Proto-Indo-European \emph{*wer-}, \textquotedblleft{}to guard, to protect\textquotedblright{} (EtymDB 2.1).
\end{cousinweb}

\begin{grammarlens}[title={What you've built}]
One more word for this chapter.
\end{grammarlens}

\subsection*{Your turn}
\begin{itemize}
\raggedright
  \item \emph{Say it:} \emph{verēcundus, verēcunda, verēcundum}
  \item \emph{Say it:} \emph{verēcundus, verēcunda, verēcundum}, once more
  \item \emph{Say it:} say \emph{verēcundus, verēcunda, verēcundum}
  \item \emph{Recall:} say the Latin for of the country, then the Latin for cheerful, then say \emph{verēcundus, verēcunda, verēcundum} again
\end{itemize}

\begin{tcolorbox}[breakable,colback=teal!4,colframe=teal!35!black,title={Before you move on}]
What does verēcundus, verēcunda, verēcundum mean? (\textquotedblleft{}Shy, modest\textquotedblright{}.) Say it once more.
\end{tcolorbox}
